%! Matrices and Column Spaces — Unit 1, Lesson 1.3 — Exit Ticket (Answer Key)
\documentclass[10pt]{article}
\usepackage{linalg-article}
\usepackage{linalg-key}   % pulls in -boxes; do NOT also load -boxes
\newcommand{\TallMath}[1]{$\displaystyle #1\rule[-1.4em]{0pt}{3.2em}$}
\newcommand{\xx}{\mathbf{x}}
\newcommand{\bb}{\mathbf{b}}

\begin{document}

\pageheader{Unit 1, Lesson 1.3}{Exit Ticket --- Answer Key}
\namedateperiod

\vspace{0.15in}

Show your work.

\begin{enumerate}[leftmargin=1.8em, itemsep=14pt]
  \item Name the two columns of $A=\begin{bmatrix} 1 & 4 \\ 2 & 3 \end{bmatrix}$, then compute
        $A\begin{bmatrix} 2 \\ 1 \end{bmatrix}$ as a combination of the columns.
        \ansline{Columns $\begin{bsmallmatrix} 1\\2 \end{bsmallmatrix}$ and $\begin{bsmallmatrix} 4\\3 \end{bsmallmatrix}$. \ $2\begin{bsmallmatrix} 1\\2 \end{bsmallmatrix}+1\begin{bsmallmatrix} 4\\3 \end{bsmallmatrix}=\begin{bsmallmatrix} 2\\4 \end{bsmallmatrix}+\begin{bsmallmatrix} 4\\3 \end{bsmallmatrix}=\begin{bsmallmatrix} 6\\7 \end{bsmallmatrix}$.}
  \item The matrix $B=\begin{bmatrix} 1 & 2 \\ 2 & 4 \end{bmatrix}$ has parallel columns. Is
        $\begin{bmatrix} 1 \\ 1 \end{bmatrix}$ in the column space $C(B)$? Decide using the columns,
        and say what your answer means about solving $B\xx=\begin{bsmallmatrix} 1\\1 \end{bsmallmatrix}$.
        \ansline{Both columns are multiples of $\begin{bsmallmatrix} 1\\2 \end{bsmallmatrix}$, so $C(B)$ is that line. $\begin{bsmallmatrix} 1\\1 \end{bsmallmatrix}$ is not a multiple of $\begin{bsmallmatrix} 1\\2 \end{bsmallmatrix}$, so it is \emph{not} in $C(B)$ --- and $B\xx=\begin{bsmallmatrix} 1\\1 \end{bsmallmatrix}$ has \emph{no solution}.}
\end{enumerate}

\begin{teachernote}
Item~2 checks the central idea: $\bb$ is reachable exactly when it lies in the column space, and
parallel columns collapse that space to a line. Full credit needs both the check (not a multiple
of $\begin{bsmallmatrix} 1\\2 \end{bsmallmatrix}$) \emph{and} the meaning (no solution). A common
slip in item~1 is crossing components --- look for $2\begin{bsmallmatrix} 1\\2 \end{bsmallmatrix}+1\begin{bsmallmatrix} 4\\3 \end{bsmallmatrix}$, not a row-by-row mash.
\end{teachernote}

\end{document}
